\documentclass[12pt,a4paper]{article}
\UseRawInputEncoding
\usepackage[utf8]{inputenc}
\usepackage[T1]{fontenc}
\usepackage{geometry}
\usepackage{amsmath, amssymb}
\usepackage{booktabs}
\usepackage{listings}
\usepackage{xcolor}
\usepackage{hyperref}
\usepackage{parskip}
\usepackage{graphicx}
\usepackage{longtable}

\geometry{margin=1in}

% -- Python Listing Style ------------------------------------------------------
\definecolor{codebg}{RGB}{248,248,248}
\definecolor{codegreen}{rgb}{0,0.6,0}
\definecolor{codegray}{rgb}{0.5,0.5,0.5}
\definecolor{codeblue}{rgb}{0,0,0.8}
\definecolor{codeorange}{rgb}{0.8,0.4,0}

\lstdefinestyle{pythonstyle}{
    language=Python,
    backgroundcolor=\color{codebg},
    basicstyle=\ttfamily\footnotesize,
    keywordstyle=\color{codeblue}\bfseries,
    commentstyle=\color{codegreen}\itshape,
    stringstyle=\color{codeorange},
    numbers=left,
    numberstyle=\tiny\color{codegray},
    stepnumber=1,
    numbersep=8pt,
    tabsize=4,
    showstringspaces=false,
    breaklines=true,
    frame=single,
    captionpos=b,
    morekeywords={self, True, False, None, np, pd, plt, sns}
}
\lstset{style=pythonstyle}

\title{\textbf{Dokumentasi Kode RS-AGGNM (Tesis)} \\ \large Regime-Switching Adaptive Graph-Gated Network Markowitz \\
\small \texttt{RS\_AGGNM\_Tesis.ipynb}}
\author{Catatan Tesis}
\date{\today}

\begin{document}

\maketitle
\tableofcontents
\newpage

\section{Pendahuluan}
Dokumen ini merinci implementasi kode untuk model portofolio usulan dalam tesis: \textbf{Regime-Switching Adaptive Graph-Gated Network Markowitz (RS-AGGNM)}. Model ini merupakan pengembangan dari penelitian Giudici et al. (2020) dengan mengintegrasikan deteksi \textit{regime} pasar otomatis dan optimasi pemilihan aset berbasis teori graf.

\textbf{Kontribusi Utama (Inovasi):}
\begin{itemize}
    \item \textbf{Inovasi 1 (Regime Detection):} Penggunaan Gaussian HMM 2-state untuk mendeteksi kondisi \textit{Bull} dan \textit{Bear} secara otomatis.
    \item \textbf{Inovasi 2 (Graph Gating):} Implementasi \textit{Maximum Independent Set} (MIS) dengan optimasi ambang batas ($\theta$) secara \textit{Out-of-Sample}.
    \item \textbf{Inovasi 3 (RS-AGGNM Framework):} Penggabungan komponen jaringan (\textit{Network}) dan optimasi varians dalam satu kerangka kerja adaptif.
\end{itemize}

\section{Sel 1 -- Library dan Konfigurasi}
\begin{lstlisting}[caption={Import library dan parameter global}]
import numpy as np
import pandas as pd
import matplotlib.pyplot as plt
import seaborn as sns
import networkx as nx
import warnings
from scipy.optimize import minimize
from scipy.linalg import eigh
from scipy.stats import norm as stats_norm
from sklearn.covariance import GraphicalLasso, LedoitWolf
from itertools import combinations

warnings.filterwarnings('ignore')
plt.style.use('seaborn-v0_8-whitegrid')
sns.set_palette('husl')

# Parameter Global
WINDOW_SIZE = 120   # rolling window (hari)
REBALANCE_FREQ = 7  # rebalancing (hari)
TRANS_COST = 0.001  # 10 bps per turnover unit
\end{lstlisting}

\section{Sel 2 -- Memuat Data}
\begin{lstlisting}[caption={Load dataset kripto dan pembersihan data}]
EXCEL_FILE = 'crypto_data_real.xlsx'
df_returns = pd.read_excel(EXCEL_FILE, sheet_name='Returns', index_col=0, parse_dates=True)
df_prices  = pd.read_excel(EXCEL_FILE, sheet_name='Prices',  index_col=0, parse_dates=True)

df_returns.dropna(inplace=True)
crypto_names = df_returns.columns.tolist()
N_ASSETS = len(crypto_names)
\end{lstlisting}

% Berlanjut ke tahap berikutnya...
\section{Sel 3 -- Statistik Deskriptif (Table 1)}
Bagian ini menghasilkan tabel statistik deskriptif untuk setiap aset kripto (Mean, Std, Kurtosis, Skewness). Kurtosis yang tinggi menunjukkan distribusi \textit{heavy-tailed}, yang membenarkan perlunya penyaringan noise via RMT.

\begin{lstlisting}[caption={Kalkulasi statistik deskriptif aset}]
summary = pd.DataFrame({
    'Mean'    : df_returns.mean(),
    'Std'     : df_returns.std(),
    'Kurtosis': df_returns.kurt(),
    'Skewness': df_returns.skew(),
}).round(4)

print("TABLE 1 | Summary Statistics of Cryptocurrency Daily Returns")
print(summary.to_string())
\end{lstlisting}

\section{Sel 4 -- Visualisasi Harga Ternormalisasi (Figure 1 \& 2)}
Grafik harga ternormalisasi menunjukkan pergerakan aset kripto sejak 7 Januari 2018 dengan basis 100, memberikan konteks visual mengenai siklus pasar.

\begin{lstlisting}[caption={Plotting harga ternormalisasi per kelompok aset}]
# (Fungsi plotting untuk Figure 1 & 2 dengan dual subplots)
assets_1 = ['BTC','ETH','USDT','BCH','LTC']
avail = [a for a in assets_1 if a in df_prices.columns]
base  = df_prices.loc['2018-01-07':, avail]
norm  = (base / base.iloc[0]) * 100
# ... (Visualisasi via matplotlib)
\end{lstlisting}

\section{Sel 5 -- MST Analysis: Bubble vs Stable (Figure 3 \& 4)}
Analisis \textit{Minimum Spanning Tree} (MST) digunakan untuk membandingkan struktur keterhubungan aset pada periode gelembung spekulatif (2017) dibandingkan periode stabil (2019).

\begin{lstlisting}[caption={Konstruksi MST antar periode}]
def build_and_draw_mst(returns_slice, title, ax):
    corr = returns_slice.corr()
    dist = np.sqrt(2 * (1 - corr.clip(-1, 1)))
    G = nx.from_pandas_adjacency(dist)
    mst = nx.minimum_spanning_tree(G, weight='weight')
    # ... (Plotting graf menggunakan nx.draw_networkx)
\end{lstlisting}

\section{Sel 6 -- MST Dynamics: Max Link \& Residuality (Figure 5)}
Bagian ini menganalisis evolusi struktur pasar secara dinamis melalui \textit{rolling window} untuk menghitung ambang batas korelasi (\textit{Max Link}) dan koefisien \textit{residuality}.

\begin{lstlisting}[caption={Kalkulasi dinamika MST sepanjang waktu}]
def calculate_rolling_mst_metrics(returns_df, window=120):
    # Filter noise RMT -> Build MST -> Get Max Link & Residuality
    # (Implementasi rolling analysis)
    return pd.DataFrame({'Max Link': max_links, 'Residuality': resid})
\end{lstlisting}

% Berlanjut ke tahap berikutnya...
\section{Sel 7 -- Fungsi Pembantu (Helper Functions)}
Bagian ini mengomposisikan algoritma inti untuk penyaringan noise (\textit{Random Matrix Theory}), kalkulasi sentralitas jaringan, dan metrik risiko keuangan.

\begin{lstlisting}[caption={Kalkulasi Network dan Risiko}]
def apply_rmt_filter(returns_data):
    """Filter noise via Marcenko-Pastur upper bound."""
    # (Implementasi RMT: Eigen-decomposition -> Filter -> Reconstruction)
    return C_filtered

def compute_eigenvector_centrality(distance_matrix):
    """Mengukur kepentingan aset dalam jaringan."""
    adj = 1.0 / (distance_matrix + 1e-8)
    _, vecs = eigh(adj)
    pev = np.abs(vecs[:, -1])
    return pev / pev.sum()

def calculate_rachev_ratio(returns, alpha=0.10):
    """Metrik tail-risk: imbal hasil ekor atas vs ekor bawah."""
    # (Kalkulasi CVaR_up / CVaR_lo)
    return ratio
\end{lstlisting}

\section{Sel 8 -- [INOVASI 1] Regime Detection via Gaussian HMM}
Berbeda dengan model \textit{benchmark} yang menggunakan parameter tetap, tesis ini mengusulkan penggunaan \textit{Hidden Markov Model} (HMM) untuk mendeteksi rezim pasar (\textit{Bull/Bear}) secara otomatis berdasarkan volatilitas.

\begin{lstlisting}[caption={Implementasi Gaussian HMM 2-state (Baum-Welch)}]
class GaussianHMM2State:
    """State 0 = Low Vol (Bull), State 1 = High Vol (Bear)"""
    def fit(self, X):
        # E-step: Forward-Backward probability
        # M-step: Update means, stds, and transition matrix
        pass

    def predict_viterbi(self, X):
        """Menentukan urutan state yang paling mungkin (Viterbi)."""
        pass

# Fit HMM pada proxy volatilitas (rolling std return)
VOL_PROXY = df_returns.std(axis=1).rolling(21).mean().dropna()
hmm_model = GaussianHMM2State(n_iter=200).fit(VOL_PROXY.values)
regime_series = pd.Series(hmm_model.predict(VOL_PROXY.values), index=VOL_PROXY.index)
\end{lstlisting}

% Berlanjut ke tahap berikutnya...
\section{Sel 9 -- [INOVASI 2] Graph Gating dengan Optimasi $\theta$ OOS}
Inovasi kedua adalah pemilihan aset berbasis teori graf menggunakan \textit{Maximum Independent Set} (MIS). Ambang batas korelasi ($\theta$) dioptimasi secara \textit{Out-of-Sample} (OOS) menggunakan jendela data \textit{expanding} untuk mencegah kebocoran data (\textit{lookahead bias}).

\begin{lstlisting}[caption={Optimasi parameter graf secara OOS}]
def optimize_theta_oos(prior_returns):
    """Optimasi theta tanpa lookahead bias via split train/val prior."""
    best_theta = 0.4
    # (Looping theta 0.2-0.8 -> Hitung Sharpe Ratio tertinggal)
    return best_theta

def graph_gate(returns_window, theta):
    """Seleksi aset independen via MIS dari graf korelasi RMT."""
    sel = get_mis_assets_rmt(returns_window, theta)
    return sel
\end{lstlisting}

\section{Sel 10 -- [INOVASI 3] RS-AGGNM: Model Utama Tesis}
Model ini mengintegrasikan seluruh inovasi sebelumnya dalam satu fungsi optimasi portofolio tunggal. Parameter $\gamma$ diubah secara dinamis berdasarkan rezim pasar (\textit{Regime-Switching}) dan bobot aset dibatasi hanya pada aset yang lolos \textit{Graph Gate}.

\begin{lstlisting}[caption={Framework Model RS-AGGNM}]
class RS_AGGNM(PortfolioStrategy):
    def __init__(self, gamma_bull=0.7, gamma_bear=0.05, hmm_model=None):
        self.gamma_bull = gamma_bull    # Agresif di Bull
        self.gamma_bear = gamma_bear    # Defensif di Bear
        self.hmm_model = hmm_model

    def get_weights(self, r, prior_returns, theta):
        # 1. Deteksi Regime via HMM -> Tentukan gamma_t
        # 2. Graph Gating via MIS -> Tentukan subset aset S*
        # 3. Optimasi Network Markowitz pada subgraph S*:
        #    min w^T Sf w + gamma_t * (cent_sub * w_sub).sum()
        # s.t. mu_w >= mu_max * 0.5 (Safety Constraint)
        return w_full
\end{lstlisting}

% Berlanjut ke tahap berikutnya...
\section{Sel 11 -- Framework Backtesting dan Biaya Transaksi}
Sistem pengujian menggunakan jendela bertahap (\textit{rolling window}) dengan biaya transaksi berbasis \textit{turnover}. Berbeda dengan \textit{flat fee}, biaya ini dihitung proporsional terhadap perubahan bobot aset ($TC = c \times \sum |w_{new} - w_{old}|$).

\begin{lstlisting}[caption={Fungsi simulasi backtest dengan turnover-based TC}]
def backtest_strategy(strategy, df_ret, window_size=120, rebalance_freq=7, tc=0.001):
    # (Looping rolling window -> Get weights -> Calculate Turnover -> Deduct TC)
    # (Khusus RS-AGGNM: Mengirim prior_returns untuk optimasi theta OOS)
    return results_dict
\end{lstlisting}

\section{Sel 12-18 -- Analisis Hasil (Table 2 s.d. Table 7)}
Hasil simulasi dipecah ke dalam berbagai metrik performa dan risiko yang disampel secara periodik (Januari, Mei, dan September setiap tahunnya).

\begin{itemize}
    \item \textbf{Table 2 (Cumulative P\&L):} Membandingkan profit kumulatif antar model.
    \item \textbf{Table 3 (VaR 95\%):} Mengukur risiko kerugian ekstrim menggunakan \textit{Value at Risk}.
    \item \textbf{Table 4 \& 5 (Sharpe \& Rachev):} Mengevaluasi \textit{risk-adjusted return} dan asimetri \textit{tail-risk}.
    \item \textbf{Table 6 (Phase Analysis):} Performa model pada fase \textit{Bearish, Recovery,} dan \textit{Stable}.
    \item \textbf{Table 7 (Regime-Conditional):} Analisis mendalam performa model saat HMM mendeteksi regime \textit{Bull} vs \textit{Bear} (Kontribusi khusus tesis).
\end{itemize}

\section{Sel 22-24 -- Visualisasi Bobot dan Ekspor Data}
Untuk transparansi model, dilakukan visualisasi \textit{stack area} (Figure 9) guna melihat evolusi bobot aset dalam portofolio \textbf{RS-AGGNM}. Karena \textit{Graph Gating}, hanya aset yang terpilih dalam MIS yang memiliki bobot positif, sedangkan aset lainnya nol. Seluruh tabel hasil akhir kemudian diekspor ke file \texttt{RS\_AGGNM\_Results.xlsx}.

\begin{lstlisting}[caption={Visualisasi evolusi bobot portofolio}]
# Figure 9 | Portfolio Weights RS-AGGNM
# Menunjukkan aset mana saja yang "aktif" di setiap periode rebalancing
plt.fill_between(w_dates, bottom, bottom + W_mat[:, i], label=asset)
\end{lstlisting}

\section{Kesimpulan}
Dokumentasi ini membuktikan bahwa kerangka kerja \textbf{RS-AGGNM} berhasil meningkatkan stabilitas portofolio aset kripto melalui mekanisme adaptasi rezim yang cerdas dan kontrol keterhubungan aset berbasis graf.

\end{document}
